\documentclass[a4paper,12pt]{article}
\usepackage[T2A]{fontenc}
\usepackage[utf8]{inputenc}
\usepackage[russian]{babel}
\usepackage{amsmath,amssymb,mathtools}
\usepackage{helvet}
\renewcommand{\familydefault}{\sfdefault}
\usepackage{icomma}
\usepackage{geometry}
\geometry{left=23mm,right=22mm,top=23mm,bottom=24mm}
\usepackage{microtype}
\usepackage{tikz}
\usetikzlibrary{arrows.meta,positioning,shapes.geometric,calc}
\usepackage{tabularx,booktabs,array}
\usepackage{enumitem}
\usepackage{hyperref}
\usepackage{parskip}
\usepackage{fancyhdr}
\usepackage{xcolor}
\usepackage{setspace}
\onehalfspacing
\definecolor{accent}{HTML}{008080}
\definecolor{darkaccent}{HTML}{00575A}
\definecolor{ink}{HTML}{292E31}
\definecolor{hint}{HTML}{596468}
\hypersetup{colorlinks=true,linkcolor=darkaccent,urlcolor=darkaccent}
\setlength{\parindent}{0pt}
\setlist[itemize]{leftmargin=6mm,itemsep=1pt,topsep=3pt}
\setlist[enumerate]{leftmargin=7mm,itemsep=2pt,topsep=3pt}
\setlength{\emergencystretch}{3em}
\renewcommand{\arraystretch}{1.28}
\pagestyle{fancy}
\fancyhf{}
\fancyhead[L]{\footnotesize\color{hint} Средняя скорость и задачи на движение}
\fancyhead[R]{\footnotesize\color{hint} Анна Баннова}
\fancyfoot[C]{\footnotesize\color{hint}\thepage}
\renewcommand{\headrulewidth}{0.25pt}
\renewcommand{\footrulewidth}{0pt}
\newcommand{\key}[1]{\textcolor{darkaccent}{\textbf{#1}}}
\newcommand{\task}[2]{\par\smallskip\noindent\textbf{№ #1. #2}\par\nobreak\smallskip}
\newcommand{\sol}{\textcolor{darkaccent}{\textbf{Решение.}}\enspace}
\newcolumntype{Y}{>{\raggedright\arraybackslash}X}
\begin{document}
\thispagestyle{empty}
\begin{center}
\vspace*{25mm}
{\Large\color{darkaccent} МАТЕМАТИКА}\par
\vspace{22mm}
{\fontsize{29}{34}\selectfont\color{ink} Чек-лист}\par
\vspace{12mm}
{\fontsize{22}{28}\selectfont\color{accent} Средняя скорость\\[3pt] и задачи на движение}\par
\vspace{19mm}
{\large Подготовка к ЕГЭ по математике}\par
\vfill
{\Large Анна Баннова}\par
\vspace{11mm}
{\large 13 ноября 2026 г.}\par
\vspace{18mm}
{\normalsize Лёвин Артём Александрович\\эксклюзивно для Анны Банновой}\par
\vspace{14mm}
\begin{tikzpicture}[x=1cm,y=1cm]
\draw[darkaccent,line width=0.9pt] (-5.2,0) .. controls (-3.5,0.9) and (-1.5,-0.7) .. (0,0) .. controls (1.5,0.7) and (3.3,-0.8) .. (5.2,0);
\fill[darkaccent] (0,0) circle (1.7pt);
\end{tikzpicture}
\end{center}
\clearpage

\section*{1. Главное правило}
Средняя скорость за весь маршрут показывает, с какой \emph{постоянной} скоростью можно было бы пройти тот же общий путь за то же общее время. Она определяется отношением \key{всего пути} к \key{всему времени}:
\[
\boxed{\displaystyle v_{\text{ср}}=\frac{S_{\text{общ}}}{t_{\text{общ}}}=\frac{S_1+S_2+\cdots+S_n}{t_1+t_2+\cdots+t_n}.}
\]
Для каждого участка пользуемся одной из трёх равносильных формул:
\[
S_i=v_it_i,\qquad t_i=\frac{S_i}{v_i},\qquad v_i=\frac{S_i}{t_i}.
\]
В задачах о движении по маршруту без остановок сумма времён равна общему времени. Если остановка включена в интервал, за который требуется найти среднюю скорость, её продолжительность также входит в знаменатель.

\key{Главная проверка формулировки:} какую величину делят на части --- \emph{время} или \emph{расстояние}? От этого зависит заполнение таблицы.

\section*{2. Две половины: принципиально разные случаи}
\begin{tabularx}{\linewidth}{@{}p{0.23\linewidth}YY@{}}
\toprule
& \textbf{Равные промежутки времени} & \textbf{Равные участки пути} \\
\midrule
В условии & «Первую половину времени ехали со скоростью $v_1$, вторую --- $v_2$». & «Первую половину пути ехали со скоростью $v_1$, вторую --- $v_2$». \\
Что известно & $t_1=t_2=T/2$ & $S_1=S_2=S/2$ \\
Что найти & $S_1=v_1T/2$, $S_2=v_2T/2$ & $t_1=S/(2v_1)$, $t_2=S/(2v_2)$ \\
Итог & $\displaystyle v_{\text{ср}}=\frac{v_1+v_2}{2}$ & $\displaystyle v_{\text{ср}}=\frac{2v_1v_2}{v_1+v_2}$ \\
\bottomrule
\end{tabularx}

\textbf{Пример 1. Равные половины времени.} Автомобиль первую половину времени двигался со скоростью $84$ км/ч, вторую половину --- $56$ км/ч. Если всё время равно $T$, то расстояния равны $42T$ и $28T$. Поэтому
\[
v_{\text{ср}}=\frac{42T+28T}{T}=\boxed{70\ \text{км/ч}}.
\]
\textbf{Вывод:} среднее арифметическое двух скоростей применимо здесь благодаря равным временам.

\clearpage
\section*{3. Равные половины пути: подробный расчёт}
\textbf{Пример 2.} Машина прошла первую половину трассы со скоростью $72$ км/ч, вторую половину --- $108$ км/ч. Найдите среднюю скорость на всём пути.

Обозначим полный путь буквой $S$, где $S>0$. Составим таблицу: 
\begin{center}
\begin{tabular}{@{}lccc@{}}
\toprule
Участок & Скорость, км/ч & Время, ч & Путь, км \\
\midrule
Первый & $72$ & $\dfrac{S}{144}$ & $\dfrac{S}{2}$ \\[8pt]
Второй & $108$ & $\dfrac{S}{216}$ & $\dfrac{S}{2}$ \\
\bottomrule
\end{tabular}
\end{center}
\noindent Время первого участка получили так: $(S/2):72=S/144$; второго --- аналогично. Подставляем данные в определение средней скорости:
\[
 v_{\text{ср}}=\frac{S/2+S/2}{S/144+S/216}
 =\frac{S}{\frac{216S+144S}{144\cdot216}}
 =\frac{S}{\frac{360S}{144\cdot216}}
 =\frac{144\cdot216}{360}.
\]
Чтобы не перемножать большие числа заранее, разложим их на простые множители:
\[
144=2^4\cdot3^2,\quad 216=2^3\cdot3^3,\quad 360=2^3\cdot3^2\cdot5.
\]
Тогда
\[
\frac{144\cdot216}{360}=
\frac{2^4\cdot3^2\cdot2^3\cdot3^3}{2^3\cdot3^2\cdot5}
=\frac{2^4\cdot3^3}{5}=\frac{16\cdot27}{5}=\boxed{86{,}4\ \text{км/ч}}.
\]
\key{Почему не подходит $(72+108)/2$?} На более медленную половину трассы затрачено больше времени. Полученное обычным усреднением число $90$ км/ч завысило бы результат.

\section*{4. Как раскрывать многоэтажные дроби}
Для $b\ne0$, $c\ne0$, $d\ne0$ верны преобразования:
\[
\frac{\frac ab}{\frac cd}=\frac{ad}{bc},\qquad
\frac{a}{\frac bc}=\frac{ac}{b},\qquad
\frac{\frac ab}{c}=\frac{a}{bc}.
\]
Общий приём: \key{деление на дробь заменяем умножением на обратную дробь}. Можно мысленно представить число $a$ как $a/1$ (так называемая «фантомная единица»). Перед сокращением проверьте, что сокращаемый множитель отличен от нуля.

\section*{5. Три участка: делим время на трети}
\textbf{Пример 3.} Первую треть всего времени автомобиль двигался со скоростью $50$ км/ч, вторую --- $75$ км/ч, последнюю --- $85$ км/ч.

Пусть суммарное время равно $T>0$. На каждый участок приходится $T/3$ часа:
\[
\begin{array}{c|c|c}
\text{Скорость}&\text{Время}&\text{Расстояние}\\ \hline
50&T/3&50T/3\\
75&T/3&75T/3\\
85&T/3&85T/3
\end{array}
\qquad
v_{\text{ср}}=\frac{50T/3+75T/3+85T/3}{T}
=\frac{210}{3}=\boxed{70\ \text{км/ч}}.
\]
Если промежутки времени \emph{равны}, средняя скорость за них равна среднему арифметическому скоростей. Это верно для любого числа равных промежутков.

\textbf{Пример 4.} Первые два часа автомобиль ехал со скоростью $55$ км/ч, следующий час --- $70$ км/ч, последние три часа --- $90$ км/ч.
\[
\begin{array}{c|c|c}
\text{Скорость}&\text{Время}&\text{Расстояние}\\ \hline
55&2&110\\70&1&70\\90&3&270
\end{array}
\qquad
v_{\text{ср}}=\frac{110+70+270}{2+1+3}
=\frac{450}{6}=\boxed{75\ \text{км/ч}}.
\]
Здесь времена различны, поэтому скорости взвешиваются по продолжительности каждого участка.

\section*{6. Проехали треть пути: сколько осталось?}
\textbf{Пример 5.} Водитель проехал треть дороги со скоростью $60$ км/ч. Оставшийся путь он преодолел со скоростью $80$ км/ч. Какова средняя скорость?

\begin{center}
\begin{tikzpicture}[x=1cm,y=1cm]
\draw[very thick,accent] (0,0) -- (4.0,0);
\draw[very thick,darkaccent] (4.0,0) -- (12,0);
\foreach \x in {0,4,12}{\draw[ink] (\x,-0.13)--(\x,0.13);}
\node[above] at (2,0.10) {$S/3$};
\node[above] at (8,0.10) {$2S/3$};
\node[below] at (6,-0.17) {Полный путь $S$};
\end{tikzpicture}
\end{center}
Оставшаяся часть: $S-S/3=2S/3$. Времена движения равны
\[
t_1=\frac{S/3}{60}=\frac{S}{180},\qquad
t_2=\frac{2S/3}{80}=\frac{S}{120}.
\]
Следовательно,
\[
v_{\text{ср}}=\frac{S}{S/180+S/120}
=\frac{1}{1/180+1/120}=\boxed{72\ \text{км/ч}}.
\]
\key{Для любой доли пути:} если проехали $\alpha S$, то осталось $(1-\alpha)S$. Например, после одной четверти осталось три четверти: $S-S/4=3S/4$.

\clearpage
\section*{7. Движение по реке}
Если собственная скорость лодки в стоячей воде равна $v$ км/ч, а скорость течения $u$ км/ч, где $v>u\ge0$, то
\[
\boxed{v_{\text{по течению}}=v+u,\qquad v_{\text{против течения}}=v-u.}
\]
\textbf{Пример 6.} Собственная скорость катера $22$ км/ч, скорость течения $2$ км/ч. По течению он двигался $6$ часов, против течения --- $2$ часа. Найдём среднюю скорость за всё время движения.
\[
\begin{array}{c|c|c|c}
\text{Направление}&\text{Скорость}&\text{Время}&\text{Путь}\\ \hline
\text{По течению}&22+2=24&6&144\\
\text{Против течения}&22-2=20&2&40
\end{array}
\]
\[
v_{\text{ср}}=\frac{144+40}{6+2}=\frac{184}{8}=\boxed{23\ \text{км/ч}}.
\]
Обратите внимание: скорость течения влияет на скорость лодки относительно берега. В общую формулу средней скорости подставляются именно \emph{фактически пройденные расстояния} и \emph{соответствующие времена}.

\section*{8. Контрольные точки перед ответом}
\begin{enumerate}[label=\arabic*.]
\item Выделите в тексте каждый отдельный участок маршрута и составьте таблицу «скорость --- время --- путь».
\item Прочитайте слова «половина пути», «треть времени», «оставшийся путь», «первые два часа». Для долей времени вводите $T$, для долей пути --- $S$.
\item Заполните таблицу при помощи $S=vt$ и $t=S/v$; при наличии реки сначала найдите $v+u$ и $v-u$.
\item Сложите \emph{все расстояния}, отдельно сложите \emph{все времена}; затем разделите первую сумму на вторую.
\item При необходимости приведите дроби к общему знаменателю, уберите многоэтажную дробь, сократите общие множители. Не сокращайте отдельные слагаемые суммы.
\item Убедитесь, что ответ имеет размерность км/ч и для движения с положительными скоростями лежит между минимальной и максимальной скоростями участков.
\end{enumerate}
\textbf{Типичные ошибки:} путают долю времени и долю пути; неверно находят оставшийся путь; меняют знак течения; усредняют скорости без учёта времени; делят сумму скоростей на сумму времён; сокращают отдельные слагаемые суммы.

\clearpage
\thispagestyle{fancy}
\begin{center}
{\Large\color{darkaccent}\textbf{Алгоритм: как найти среднюю скорость}}\par
\vspace{5mm}
{\small Сначала извлеките из условия данные, затем выполняйте вычисления.}
\end{center}
\vspace{10mm}
\begin{center}
\begin{tikzpicture}[
  node distance=15mm and 13mm,
  action/.style={draw=accent,rounded corners=2pt,very thick,align=center,text width=105mm,minimum height=13mm,inner sep=5pt,fill=white},
  decision/.style={diamond,draw=darkaccent,very thick,aspect=2.6,align=center,text width=49mm,inner sep=1.5pt,fill=white},
  branch/.style={draw=accent,rounded corners=2pt,align=center,text width=57mm,minimum height=20mm,inner sep=4pt,fill=white},
  line/.style={-{Stealth[length=3mm]},thick,darkaccent},
  every node/.style={font=\small,inner sep=3pt}
]
\node[action] (start) {1. Разделите весь маршрут на участки.\\ Укажите скорость на каждом участке.};
\node[action,below=of start] (river) {2. Для реки: по течению $v+u$, против течения $v-u$.};
\node[decision,below=17mm of river] (which) {Что дано по участкам?};
\node[branch] (times) at ([xshift=-41mm,yshift=-43mm]which.south) {\textbf{Время или его доля}\\ Найдите путь: $S_i=v_i t_i$.};
\node[branch] (paths) at ([xshift=41mm,yshift=-43mm]which.south) {\textbf{Путь или его доля}\\ Найдите время: $t_i=S_i/v_i$.};
\node[action] (sum) at ([yshift=-80mm]which.south) {3. Посчитайте $S_{\text{общ}}=\sum S_i$ и $t_{\text{общ}}=\sum t_i$.};
\node[action,below=of sum] (finish) {4. Вычислите $v_{\text{ср}}=S_{\text{общ}}/t_{\text{общ}}$.\\ Проверьте единицы измерения и разумность ответа.};
\coordinate (join) at ([yshift=12mm]sum.north);
\draw[line] (start) -- (river);
\draw[line] (river) -- (which);
\draw[line] (which.west) -| (times.north);
\draw[line] (which.east) -| (paths.north);
\draw[thick,darkaccent] (times.south) |- (join);
\draw[thick,darkaccent] (paths.south) |- (join);
\draw[line] (join) -- (sum);
\draw[line] (sum) -- (finish);
\end{tikzpicture}
\end{center}
\clearpage

\section*{Тренировочный блок}
\begingroup\setstretch{1.30}
\textcolor{hint}{Путь А: 10 задач. Условия расположены по нарастанию сложности.}\par
\smallskip
Во всех задачах требуется найти \key{среднюю скорость на всём указанном маршруте}. Если специально не оговорено иное, транспорт двигался без остановок, а указанные скорости постоянны на соответствующих участках. Составляйте таблицу «скорость --- время --- расстояние» и записывайте формулу для средней скорости перед вычислением.

\task{1}{Доставка по загородной дороге}
Курьер выехал из небольшого посёлка в районный центр. Первые $2$ часа он двигался по просёлочной дороге с постоянной скоростью $48$ км/ч. После выезда на шоссе курьер ещё $3$ часа ехал со скоростью $64$ км/ч. Остановок на маршруте не было. Определите среднюю скорость курьера от отправления из посёлка до прибытия в районный центр.

\task{2}{Две равные части времени}
Водитель возвращался домой по дороге, где участки различались по разрешённой скорости. В течение первой половины всего времени поездки он ехал со скоростью $54$ км/ч, а на протяжении второй половины времени --- со скоростью $90$ км/ч. Общее время путешествия неизвестно, но обе его половины были одинаковой продолжительности. Найдите среднюю скорость автомобиля за поездку.

\task{3}{Три равные части времени}
Микроавтобус двигался по маршруту, проходящему через городскую улицу, пригородную дорогу и скоростной участок трассы. На каждую из этих трёх частей пришлось ровно по одной трети общего времени движения. Скорости на последовательных участках составляли $48$, $72$ и $96$ км/ч. Определите среднюю скорость микроавтобуса за весь маршрут.

\task{4}{Две равные части расстояния}
Автомобиль преодолел дорогу между двумя населёнными пунктами, причём её протяжённость неизвестна. Ровно первую половину расстояния водитель прошёл со скоростью $60$ км/ч, а оставшуюся половину --- со скоростью $90$ км/ч. Время движения на этих участках оказалось разным. Рассчитайте среднюю скорость на всём пути и объясните в записи решения, почему следует складывать времена.

\clearpage
\section*{Тренировочный блок. Продолжение}
\task{5}{Треть пути и оставшееся расстояние}
Легковой автомобиль направился из города в загородный туристический комплекс. Первую треть полного пути он двигался по местной дороге со скоростью $45$ км/ч. На оставшейся части маршрута, проходившей по магистрали, его скорость составила $90$ км/ч. Общая длина поездки неизвестна. Определите долю оставшегося пути, составьте таблицу для двух участков и найдите среднюю скорость.

\task{6}{Одна четверть дороги}
Во время поездки на соревнования автобус первую четверть расстояния ехал по загруженной городской дороге со скоростью $60$ км/ч. Затем он выехал на свободную трассу и преодолел всё оставшееся расстояние со скоростью $90$ км/ч. Автобус не останавливался. Обозначьте весь путь через $S$, выразите длины обоих участков через $S$ и вычислите среднюю скорость движения.

\task{7}{Участки разной продолжительности}
Грузовой автомобиль вёз оборудование между складами. Сначала он ехал $2$ часа со скоростью $50$ км/ч, затем $1$ час со скоростью $80$ км/ч и, наконец, ещё $3$ часа по скоростной дороге со скоростью $100$ км/ч. Учитывается только время движения. Найдите расстояние, пройденное на каждом участке, и среднюю скорость за всю шестичасовую поездку.

\clearpage
\section*{Тренировочный блок. Завершение}
\task{8}{Три равные части пути}
Туристический автобус пересекал три дорожные зоны, каждая из которых составляла ровно треть общей протяжённости маршрута. В первой зоне его скорость была $36$ км/ч, во второй --- $54$ км/ч, в третьей --- $108$ км/ч. Продолжительность движения на этих участках различалась. Обозначьте общий путь буквой $S$, выразите все три времени через $S$ и найдите среднюю скорость за весь маршрут.

\task{9}{По течению и против течения}
Катер выполнял рейс по реке. Его собственная скорость в стоячей воде равнялась $18$ км/ч, а скорость речного течения составляла $2$ км/ч. Первые $5$ часов катер двигался по течению, после чего развернулся и ещё $3$ часа шёл против течения. Течение было постоянным, остановок не было. Найдите скорость катера относительно берега на каждом участке, расстояния и среднюю скорость за $8$ часов движения.

\task{10}{Сложный маршрут с тремя долями пути}
Автомобиль совершил междугороднюю поездку по трём дорогам. Первую четверть всего расстояния водитель двигался со скоростью $60$ км/ч. Следующую четверть расстояния, проходившую через холмистую местность, он преодолел со скоростью $80$ км/ч. Оставшуюся половину пути автомобиль прошёл по автомагистрали со скоростью $120$ км/ч. Общая длина маршрута неизвестна. Составьте таблицу из трёх строк, найдите среднюю скорость в виде точной дроби и округлите её до десятых км/ч.
\endgroup

\clearpage
\section*{Ответы к тренировочному блоку}
\begin{center}
\begin{tabularx}{\linewidth}{@{}p{12mm}Yp{55mm}@{}}
\toprule
\textbf{№} & \textbf{Ключевая запись для проверки} & \textbf{Ответ} \\
\midrule
1 & $\displaystyle\frac{48\cdot2+64\cdot3}{2+3}$ & $57{,}6$ км/ч \\
2 & $\displaystyle\frac{54+90}{2}$ & $72$ км/ч \\
3 & $\displaystyle\frac{48+72+96}{3}$ & $72$ км/ч \\
4 & $\displaystyle\frac{2\cdot60\cdot90}{60+90}$ & $72$ км/ч \\
5 & $\displaystyle\frac{1}{1/135+2/270}$ & $67{,}5$ км/ч \\
6 & $\displaystyle\frac{1}{1/240+3/360}$ & $80$ км/ч \\
7 & $\displaystyle\frac{50\cdot2+80\cdot1+100\cdot3}{6}$ & $80$ км/ч \\
8 & $\displaystyle\frac{3}{1/36+1/54+1/108}$ & $54$ км/ч \\
9 & $\displaystyle\frac{20\cdot5+16\cdot3}{5+3}$ & $18{,}5$ км/ч \\
10 & $\displaystyle\frac{1}{1/240+1/320+1/240}$ & $\displaystyle\frac{960}{11}\approx87{,}3$ км/ч \\
\bottomrule
\end{tabularx}
\end{center}
\medskip
\textbf{Дополнительная проверка к № 5:} оставшаяся доля пути --- $2/3$. \textbf{К № 6:} осталось $3/4$ пути. \textbf{К № 8:} каждый участок равен $S/3$. \textbf{К № 9:} скорости по течению и против него равны $20$ и $16$ км/ч; пути --- $100$ и $48$ км. \textbf{К № 10:} доли пути равны $S/4$, $S/4$ и $S/2$.

\vfill
\begin{center}
\textcolor{hint}{Лёвин Артём Александрович --- эксклюзивно для Анны Банновой}\par
\smallskip
\begin{tikzpicture}
\draw[accent,line width=.7pt] (-4,0) .. controls (-2.6,.4) and (-1.4,-.3) .. (0,0) .. controls (1.4,.3) and (2.5,-.4) .. (4,0);
\end{tikzpicture}
\end{center}
\end{document}
